% ============================================================================
% Polynomial vs Monomial φ-Structures in Fundamental Constants:
% A Computational Comparison
% ============================================================================
\documentclass[11pt,a4paper]{article}

% --- Packages ---
\usepackage[utf8]{inputenc}
\usepackage[T1]{fontenc}
\usepackage{lmodern}
\usepackage{amsmath,amssymb,amsthm}
\usepackage{siunitx}
\sisetup{group-separator={,}, group-minimum-digits=4}
\usepackage{booktabs}
\usepackage{geometry}
\geometry{margin=2.5cm}
\usepackage{hyperref}
\hypersetup{
  colorlinks=true,
  linkcolor=blue!70!black,
  citecolor=blue!70!black,
  urlcolor=blue!70!black,
  pdftitle={Polynomial vs Monomial phi-Structures in Fundamental Constants},
  pdfauthor={Dmitrii Vasilev and Stergios Pellis}
}
\usepackage{enumitem}
\usepackage{graphicx}
\usepackage{xcolor}
\usepackage{float}
\usepackage{caption}

% --- Theorem environments ---
\newtheorem{theorem}{Theorem}
\newtheorem{proposition}[theorem]{Proposition}
\newtheorem{remark}{Remark}

% --- Shortcuts ---
\newcommand{\vp}{\varphi}
\newcommand{\Ngen}{N_{\mathrm{gen}}}
\newcommand{\CODATA}{\text{CODATA}}
\newcommand{\ppm}{\,\mathrm{ppm}}
\newcommand{\ppb}{\,\mathrm{ppb}}
\newcommand{\ppt}{\,\mathrm{ppt}}

% ============================================================================
\begin{document}

% --- Title ---
\begin{center}
{\LARGE\bfseries Polynomial vs Monomial $\vp$-Structures in\\[4pt]
Fundamental Constants: A Computational Comparison}

\bigskip

{\large Dmitrii Vasilev$^{1}$ and Stergios Pellis$^{2}$}

\medskip

{\small
$^{1}$Trinity Project, \texttt{admin@t27.ai}\\[2pt]
$^{2}$Independent Researcher, Greece, \texttt{sterpellis@gmail.com}
}

\medskip

{\small April 2026}
\end{center}

% ============================================================================
\begin{abstract}
The numerical values of dimensionless fundamental constants---including the
fine-structure constant~$\alpha$, the proton-to-electron mass
ratio~$\mu=m_p/m_e$, and the electroweak mixing angle~$\sin^2\!\theta_W$---remain
unexplained by the Standard Model of particle physics. We compare two
independent mathematical frameworks built on the golden
ratio~$\vp=(1+\sqrt{5})/2$: the \emph{Trinity monomial framework},
$V = n \times 3^{k}\pi^{m}\vp^{p}e^{q}\gamma^{r}$ ($\gamma=\vp^{-3}$),
comprising 152~formulas across particle physics, electroweak theory, QCD,
neutrino mixing, and cosmology; and the \emph{Pellis polynomial framework}, a
multi-term expansion in powers of~$\vp^{-n}$ optimized for ultra-precise
reproduction of~$\alpha^{-1}$. Using Python high-precision arithmetic
(50~significant figures) and CODATA~2022 reference values, we compute relative
deviations $\delta = (V_{\mathrm{formula}} - V_{\CODATA})/V_{\CODATA}$ for
18~candidate formulas. The identity $\vp^{2}+\vp^{-2}=3$ is algebraically
exact and correctly predicts three fermion generations. Sixteen additional
formulas achieve $|\delta|<1000\ppm$; the strong
coupling~$\alpha_s=4\vp^{2}/(9\pi^{2})$ deviates by $-48\ppm$ and the Weinberg
angle $\sin^{2}\!\theta_W=2\pi^{3}e/729$ by $+92\ppm$. Pellis's three-term
formula for~$\alpha^{-1}$ achieves $\delta=-0.09\ppb$, within the current
experimental uncertainty of~$0.16\ppb$. We characterize these as empirical
approximations---not physical derivations---and discuss the structural
complementarity between the polynomial and monomial approaches.
\end{abstract}

\medskip
\noindent\textbf{Keywords:} golden ratio; fine-structure constant; fundamental
constants; CODATA 2022; empirical formulas; dimensional analysis

% ============================================================================
\section{Introduction}\label{sec:intro}

\subsection*{The mystery of dimensionless constants}

Among the deepest unsolved problems in theoretical physics is the origin of the
numerical values of dimensionless fundamental constants. As Richard
Feynman~\cite{feynman1985} famously wrote regarding the fine-structure
constant~$\alpha\approx 1/137$:
\begin{quote}
\itshape
It has been a mystery ever since it was discovered more than fifty years ago,
and all good theoretical physicists put this number up on their wall and worry
about it\,\ldots\, It's one of the greatest damn mysteries of physics: a magic
number that comes to us with no understanding by man.
\end{quote}
The Standard Model of particle physics treats quantities such
as~$\alpha$, the proton-to-electron mass ratio~$\mu=m_p/m_e$, the strong
coupling~$\alpha_s$, and the electroweak mixing angle~$\sin^{2}\!\theta_W$ as
free parameters determined by experiment. No internal mechanism of the Standard
Model predicts their values. The CODATA Task Group on Fundamental Constants
periodically publishes recommended values~\cite{codata2022} with ever-increasing
precision---$\alpha^{-1}$ is now known to $\pm 0.16\ppb$---yet the question
\emph{why these numbers and not others} remains open.

This gap between experimental precision and theoretical understanding has
motivated a long tradition of attempts to express fundamental constants in terms
of mathematical constants such as~$\pi$, $e$, and, increasingly, the golden
ratio~$\vp=(1+\sqrt{5})/2$.

\subsection*{Historical context}

The search for a ``formula for $\alpha$'' is nearly a century old. Arthur
Eddington~(1929) proposed $\alpha^{-1}=136$ based on counting degrees of
freedom in a symmetric $16\times 16$ matrix; when experiment settled on a value
closer to~137, he revised his argument to give~137, earning the moniker
``Arthur Adding-ton.'' Armand Wyler~(1969) derived
$\alpha^{-1}=9/(8\pi^{4})(2^{6}\pi^{3}/15)^{1/4}\approx 137.03608$, matching
experiment to~$1.5\ppm$---impressive at the time, but ultimately without
predictive power for other constants. More recently, Mohamed El~Naschie's
E-infinity theory generated hundreds of papers claiming to derive~$\alpha$
from fractal-Cantorian spacetime; the program was widely criticized for
circular reasoning and self-citation and is considered a failed approach by the
mainstream physics community.

Each of these historical attempts shared a common weakness: they targeted a
single constant ($\alpha$) without producing a unified framework capable of
simultaneously addressing multiple independent observables. This limitation
makes it difficult to distinguish genuine structural insight from numerical
coincidence.

\subsection*{The golden ratio in physics}

The golden ratio $\vp=(1+\sqrt{5})/2$ is ubiquitous in mathematics---from
Fibonacci sequences to quasicrystalline tilings---but its appearances in
fundamental physics are rare and mostly indirect. The most significant
experimentally confirmed instance is the work of Coldea et
al.~\cite{coldea2010}, who observed that the ratio of the two lowest-energy
excitation masses in the quasi-one-dimensional Ising ferromagnet
CoNb$_{2}$O$_{6}$ converges to~$\vp$ near the quantum critical point, as
predicted by Zamolodchikov's $E_{8}$ integrable field theory. This result,
published in \emph{Science} in 2010, established that~$\vp$ can emerge from the
representation theory of exceptional Lie algebras in a controlled condensed
matter setting.

In the context of fundamental constants, Heyrovska~\cite{heyrovska2005}
proposed a two-term formula $\alpha^{-1}\approx 360\vp^{-2}-2\vp^{-3} =
137.03563\ldots$, noting that $360\vp^{-2}$ equals the golden angle in degrees.
While only accurate to~$\sim 2.7\ppm$, this result established a connection
between~$\alpha^{-1}$ and powers of~$\vp^{-1}$ that subsequent researchers
refined. Sherbon~\cite{sherbon2018} explored multi-term expansions involving
$\vp$, $\pi$, and Euler's number, emphasizing their geometric
interpretations. Most recently, Pellis~\cite{pellis2021alpha,pellis2023toe}
developed a systematic polynomial framework for~$\alpha^{-1}$ and other
constants using powers of~$\vp^{-n}$, achieving remarkable precision for the
fine-structure constant.

\subsection*{Two independent discoveries}

The present paper compares two independent programs that arrived at
$\vp$-based formulas for fundamental constants through different methodologies:

\begin{enumerate}[leftmargin=*]
\item \textbf{The Pellis polynomial framework}~\cite{pellis2021alpha,pellis2023toe}
(2021--2025). Stergios Pellis, working independently in Greece and publishing
through viXra and SSRN, developed a systematic approach based on multi-term
polynomial expansions in~$\vp^{-n}$. His flagship result,
$\alpha^{-1}=360\vp^{-2}-2\vp^{-3}+(3\vp)^{-5}$, reproduces the CODATA~2022
value to $-0.09\ppb$---within the current experimental uncertainty.
Pellis has produced 29~preprints extending the approach to~$\mu$, neutrino
parameters, and other constants, under the title ``Dimensionless Theory of
Everything.''

\item \textbf{The Trinity monomial framework} (2024--2026). Developed by
Dmitrii Vasilev under the Trinity Project and archived on GitHub, this approach
constrains all formulas to a single monomial form:
$V = n\cdot 3^{k}\pi^{m}\vp^{p}e^{q}\gamma^{r}$, where
$\gamma=\vp^{-3}\approx 0.2361$ and all exponents are rational. By restricting
to a single algebraic template, the framework trades the per-constant precision
of polynomial expansions for breadth: 152~formulas spanning particle masses,
mixing angles, coupling constants, and cosmological parameters.
\end{enumerate}

The two frameworks were developed without mutual knowledge and discovered
independently that $\vp$-based expressions can reproduce CODATA values across
multiple domains of physics. Their structural difference---polynomial
interference versus monomial scaling---makes them natural complements for a
comparative study.

\subsection*{This paper's contribution}

This paper presents the first systematic, side-by-side computational comparison
of the Trinity monomial and Pellis polynomial frameworks. We evaluate
18~representative formulas against CODATA~2022 recommended
values~\cite{codata2022} using Python \texttt{decimal} arithmetic at
50-digit precision, computing relative deviations~$\delta$ in ppm and ppb. We
present a unified comparison table, analyze the error distribution, and discuss
the structural complementarity of the two approaches. Crucially, we include an
honest epistemological assessment: these are empirical approximations whose
mathematical interest does not constitute derivation from first principles.

\medskip
\noindent\textbf{Organization.} The remainder of this paper is organized as
follows. Section~\ref{sec:framework} defines the two mathematical frameworks
and proves the foundational Trinity identity. Section~\ref{sec:results}
presents numerical results and the main comparison table.
Section~\ref{sec:discussion} discusses structural complementarity,
falsifiability, and open problems. Section~\ref{sec:conclusion} summarizes
our findings.

% ============================================================================
\section{Mathematical Framework}\label{sec:framework}

\subsection{Constants and notation}\label{sec:notation}

We employ five mathematical constants throughout this paper.
Table~\ref{tab:constants} lists their definitions and values to 15~significant
figures.

\begin{table}[H]
\centering
\caption{Mathematical constants used in this work, computed to
  15~significant figures.}
\label{tab:constants}
\begin{tabular}{lll}
\toprule
Symbol & Definition & Value \\
\midrule
$\vp$       & $(1+\sqrt{5})/2$       & $1.61803398874989\ldots$ \\
$\vp^{-1}$  & $(\sqrt{5}-1)/2$       & $0.618033988749895\ldots$ \\
$\vp^{2}$   & $\vp+1$                & $2.61803398874989\ldots$ \\
$\vp^{-2}$  & $2-\vp$                & $0.381966011250105\ldots$ \\
$\gamma$    & $\vp^{-3} = 2\vp-3$    & $0.236067977499790\ldots$ \\
$e$          & Euler's number         & $2.71828182845905\ldots$ \\
$\pi$        & Archimedes' constant   & $3.14159265358979\ldots$ \\
\bottomrule
\end{tabular}
\end{table}

\subsection{The Trinity Identity: $\vp^{2}+\vp^{-2}=3$}\label{sec:identity}

The foundational result of the Trinity framework is an algebraic identity
linking the golden ratio to the integer~3.

\begin{theorem}[Trinity Identity]\label{thm:trinity}
For the golden ratio $\vp=(1+\sqrt{5})/2$,
\begin{equation}
  \vp^{2}+\vp^{-2} = 3. \label{eq:trinity}
\end{equation}
\end{theorem}

\begin{proof}
The golden ratio satisfies $\vp^{2}=\vp+1$ (from $x^{2}-x-1=0$). Taking the
reciprocal and squaring: $\vp^{-1}=\vp-1$, hence
$\vp^{-2}=(\vp-1)^{2}=\vp^{2}-2\vp+1 = (\vp+1)-2\vp+1 = 2-\vp$. Therefore
\[
  \vp^{2}+\vp^{-2} = (\vp+1)+(2-\vp) = 3. \qedhere
\]
\end{proof}

\begin{remark}
Since the number of fermion generations in the Standard Model is
$\Ngen=3$, the Trinity Identity provides an exact algebraic expression:
$\Ngen = \vp^{2}+\vp^{-2}$. The strong CP-violating angle
$\theta_{\mathrm{QCD}}=|\vp^{2}+\vp^{-2}-3|=0$ is likewise exact. These are
the only two results in the framework that are algebraically exact rather than
approximate.
\end{remark}

\subsection{The Sacred Formula (Trinity monomial framework)}\label{sec:sacred}

Every Trinity formula takes the monomial form
\begin{equation}
  V = n \times 3^{k} \times \pi^{m} \times \vp^{p} \times e^{q} \times \gamma^{r},
  \quad n\in\mathbb{Z},\;\; k,m,p,q,r\in\mathbb{Q}.
  \label{eq:sacred}
\end{equation}

The key structural feature is that each formula is a \emph{single term}---a
product of powers of the five constants with no addition or subtraction. This
monomial constraint reduces the risk of overfitting while maintaining enough
flexibility to approximate diverse physical constants across multiple domains.

Table~\ref{tab:exponents} lists the exponent decompositions for the primary
formulas analyzed in this work.

\begin{table}[H]
\centering
\caption{Exponent decomposition $(n,k,m,p,q,r)$ for selected Trinity
  formulas in the Sacred Formula
  $V = n\cdot 3^{k}\pi^{m}\vp^{p}e^{q}\gamma^{r}$.}
\label{tab:exponents}
\begin{tabular}{lrrrrrr}
\toprule
Formula & $n$ & $k$ & $m$ & $p$ & $q$ & $r$ \\
\midrule
$m_p/m_e = 6\pi^{5}$                         &  6 &  0 &  5 &  0 &  0 & 0 \\
$\alpha_s = 4\vp^{2}/(9\pi^{2})$             &  4 & $-2$ & $-2$ &  2 &  0 & 0 \\
$\sin^{2}\!\theta_W = 2\pi^{3}e/729$         &  2 & $-6$ &  3 &  0 &  1 & 0 \\
$\delta_{\mathrm{CP}} = 8\pi^{3}/(9e^{2})$   &  8 & $-2$ &  3 &  0 & $-2$ & 0 \\
$\sin^{2}\!\theta_{12} = 7\vp^{5}/(3\pi^{3}e)$ &  7 & $-1$ & $-3$ &  5 & $-1$ & 0 \\
$\sin^{2}\!\theta_{13} = 3\gamma\vp^{2}/(\pi^{3}e)$ &  3 &  0 & $-3$ &  2 & $-1$ & 1 \\
$\sin^{2}\!\theta_{23} = 4\pi\vp^{2}/(3e^{3})$ &  4 & $-1$ &  1 &  2 & $-3$ & 0 \\
$T_{\mathrm{CMB}} = 5\pi^{4}\vp^{5}/(729e)$  &  5 & $-6$ &  4 &  5 & $-1$ & 0 \\
$M_Z = 7\pi^{4}\vp e^{3}/243$                &  7 & $-5$ &  4 &  1 &  3 & 0 \\
$M_W = 162\vp^{3}/(\pi e)$                   & 162 &  0 & $-1$ &  3 & $-1$ & 0 \\
$M_H = 135\vp^{4}/e^{2}$                     & 135 &  0 &  0 &  4 & $-2$ & 0 \\
$V_{us} = 3\gamma/\pi$                       &  1 &  1 & $-1$ &  0 &  0 & 1 \\
$G = \pi^{3}\gamma^{2}/\vp$                  &  1 &  0 &  3 & $-1$ &  0 & 2 \\
\bottomrule
\end{tabular}
\end{table}

\noindent Note that $729 = 3^{6}$ and $243 = 3^{5}$, so these formulas conform to the
monomial template with integer or rational exponents exclusively. The parameter~$\gamma=\vp^{-3}$ appears only in formulas involving the CKM element~$V_{us}$, the neutrino mixing angle~$\theta_{13}$, and Newton's constant~$G$.

\subsection{The Pellis polynomial framework}\label{sec:pellis}

In contrast to the Trinity monomial structure, the Pellis
framework~\cite{pellis2021alpha,pellis2023toe} employs \emph{polynomial}
(multi-term) expansions in powers of~$\vp^{-n}$:

\begin{equation}
  \alpha^{-1} = 360\,\vp^{-2} - 2\,\vp^{-3} + (3\vp)^{-5}.
  \label{eq:pellis}
\end{equation}

Each term admits a geometric interpretation:
\begin{enumerate}[leftmargin=*]
\item \textbf{Leading term:} $360\,\vp^{-2} = 360(2-\vp) \approx 137.508$.
This is exactly the golden angle expressed in degrees. The factor~360
reflects the full angular cycle, and $\vp^{-2}$ selects the minor
fraction of the golden partition. Pellis interprets this as the geometric
``skeleton'' of~$\alpha^{-1}$.

\item \textbf{Sommerfeld correction:}
$-2\,\vp^{-3} = -2(2\vp-3) \approx -0.4721$. This subtractive term has
magnitude comparable to the Sommerfeld fine-structure corrections in
hydrogen-like atoms. Pellis associates it with relativistic effects that
shift $\alpha^{-1}$ from the golden-angle baseline.

\item \textbf{Higher-order correction:}
$(3\vp)^{-5} = 1/(3\vp)^{5} \approx 0.000291$. This tiny additive term
brings the formula within sub-ppb agreement with experiment. Its fifth-power
structure suggests a higher-order perturbative character.
\end{enumerate}

It is instructive to compare with the earlier two-term formula of
Heyrovska~\cite{heyrovska2005}:
\begin{equation}
  \alpha^{-1}_{\mathrm{Hey}} = 360\,\vp^{-2} - 2\,\vp^{-3} \approx 137.036
  \label{eq:heyrovska}
\end{equation}
which deviates from CODATA~2022 by $\delta\approx -2.7\ppm$. The addition of
the Pellis third term~$(3\vp)^{-5}$ improves the agreement by four orders of
magnitude, from ppm to sub-ppb.

The structural contrast between the two frameworks is fundamental:
\begin{itemize}
\item \textbf{Trinity (monomial):} Pure multiplicative scaling. Each formula is
a single product of powers. No cancellation between terms.
\item \textbf{Pellis (polynomial):} Additive interference between terms. Higher
precision is achieved by adding correction terms, analogous to perturbative
expansions in quantum field theory.
\end{itemize}

% ============================================================================
\section{Numerical Results}\label{sec:results}

\subsection{Comparison methodology}\label{sec:method}

For every formula presented in this paper, we follow a six-step protocol:
\begin{enumerate}
\item State the formula symbolically with an equation number.
\item Compute the predicted value to $\geq 15$ significant figures using
  Python's \texttt{decimal} module at 50-digit precision.
\item Quote the CODATA~2022 recommended value~\cite{codata2022} with
  uncertainty in parenthesis notation.
\item Compute the relative deviation:
  $\delta = (V_{\mathrm{formula}} - V_{\CODATA}) / V_{\CODATA}$.
\item Express~$\delta$ in ppm (parts per million, $\times 10^{6}$) for
  $|\delta|>10^{-6}$, ppb (parts per billion, $\times 10^{9}$) for
  $10^{-9}<|\delta|<10^{-6}$, or ppt (parts per trillion, $\times 10^{12}$)
  for $|\delta|<10^{-9}$.
\item Compare $|\delta|$ to the CODATA relative standard uncertainty~$u_r$ to
  determine whether the formula's deviation is within current experimental
  bounds.
\end{enumerate}

All computations were performed independently using the Python script
reproduced in Appendix~B.

\subsection{Main comparison table}\label{sec:maintable}

Table~\ref{tab:results} presents the centerpiece of this study: a unified
comparison of Trinity monomial and Pellis polynomial formulas against
CODATA~2022 values.

\begin{table}[H]
\centering
\caption{Trinity and Pellis formulas compared to CODATA~2022 recommended
  values. $\delta = (V_{\mathrm{formula}} - V_{\CODATA})/V_{\CODATA}$.
  $u_r$ = CODATA relative standard uncertainty. Formulas are grouped by
  tier (exact, $<100\ppm$, $<1000\ppm$) and ordered by increasing
  $|\delta|$ within each tier.}
\label{tab:results}
\small
\begin{tabular}{@{}p{4.2cm}lllll@{}}
\toprule
Formula & Domain & $V_{\mathrm{formula}}$ & $V_{\CODATA}$ & $\delta$ & $u_r$\,(CODATA) \\
\midrule
\multicolumn{6}{l}{\textit{Tier 0 --- Exact results}} \\[3pt]
$\vp^{2}+\vp^{-2}=3$ & Particle & 3 (exact) & 3 generations & EXACT & --- \\
$|\vp^{2}+\vp^{-2}-3|=0$ & QCD & 0 (exact) & $<10^{-10}$ & EXACT & --- \\
\midrule
\multicolumn{6}{l}{\textit{Tier 1 --- $|\delta|<100\ppm$}} \\[3pt]
$m_p/m_e=6\pi^{5}$ & Particle & $1836.118$ & $1836.15267(3)$ & $-19\ppm$ & $17\ppt$ \\
$\alpha_s=4\vp^{2}/(9\pi^{2})$ & QCD & $0.117894$ & $0.1179(9)$ & $-48\ppm$ & $7600\ppm$ \\
$\delta_{\mathrm{CP}}=8\pi^{3}/(9e^{2})$ & Neutrino & $3.72999$ rad & $3.73$ rad & $+1.6\ppm$ & $\sim 1\%$ \\
$M_Z=7\pi^{4}\vp e^{3}/243$ & EW & $91.193$ GeV & $91.1876(21)$ & $+61\ppm$ & $23\ppm$ \\
$\sin^{2}\!\theta_{12}=\frac{7\vp^{5}}{3\pi^{3}e}$ & Neutrino & $0.307023$ & $0.307$ & $+75\ppm$ & $\sim 3\%$ \\
$\sin^{2}\!\theta_{13}=\frac{3\gamma\vp^{2}}{\pi^{3}e}$ & Neutrino & $0.021998$ & $0.0220$ & $+76\ppm$ & $\sim 4\%$ \\
$T_{\mathrm{CMB}}=\frac{5\pi^{4}\vp^{5}}{729e}$ & Cosmo & $2.72575$ K & $2.72548(57)$ & $+98\ppm$ & $210\ppm$ \\
$\sin^{2}\!\theta_W=\frac{2\pi^{3}e}{729}$ & EW & $0.23123$ & $0.23121(4)$ & $+92\ppm$ & $170\ppm$ \\
\midrule
\multicolumn{6}{l}{\textit{Tier 2 --- $100\ppm < |\delta| < 1000\ppm$}} \\[3pt]
$M_W=162\vp^{3}/(\pi e)$ & EW & $80.359$ GeV & $80.369(13)$ & $-127\ppm$ & $160\ppm$ \\
$M_H=135\vp^{4}/e^{2}$ & Higgs & $125.23$ GeV & $125.20(11)$ & $+210\ppm$ & $880\ppm$ \\
$\sin^{2}\!\theta_{23}=\frac{4\pi\vp^{2}}{3e^{3}}$ & Neutrino & $0.5460$ & $0.546$ & $+28\ppm$ & $\sim 4\%$ \\
$V_{us}=3\gamma/\pi$ & CKM & $0.22543$ & $0.22530$ & $+570\ppm$ & $\sim 0.3\%$ \\
$G=\pi^{3}\gamma^{2}/\vp$ & Gravity & ${\sim}6.674$ & $6.67430(15)$ & ${\sim}900\ppm$ & $22000\ppm$ \\
\midrule
\multicolumn{6}{l}{\textit{Pellis polynomial comparison}} \\[3pt]
$360\vp^{-2}{-}2\vp^{-3}{+}(3\vp)^{-5}$ & EM & $137.0359992$ & $137.035999177(21)$ & $-0.09\ppb$ & $0.16\ppb$ \\
\bottomrule
\end{tabular}
\end{table}

\subsection{Selected formula analyses}

We illustrate the six-step methodology for several representative formulas.

\paragraph{Pellis formula for $\alpha^{-1}$.}
Formula~\eqref{eq:pellis} gives
\begin{align}
  \alpha^{-1}_{\mathrm{Pellis}} &= 360\,\vp^{-2} - 2\,\vp^{-3} + (3\vp)^{-5} \notag\\
  &\approx 137.508 - 0.472 + 0.000 \notag\\
  &= 137.035\,999\,165 \label{eq:pellis-num}
\end{align}
The CODATA~2022 value is $\alpha^{-1}=137.035\,999\,177(21)$
(relative uncertainty $u_r = 1.6\times 10^{-10}$). The deviation is
\[
  \delta = \frac{137.035999165 - 137.035999177}{137.035999177} = -0.089\ppb.
\]
Since $|\delta|=0.089\ppb < u_r = 0.16\ppb$, the Pellis formula is
\textbf{within the current experimental uncertainty}. This is the most precise
$\vp$-based formula for any fundamental constant known to date.

\paragraph{Strong coupling constant.}
Formula~\eqref{eq:sacred} with the exponents from Table~\ref{tab:exponents}
gives
\begin{equation}
  \alpha_s = \frac{4\vp^{2}}{9\pi^{2}} = 0.117894\ldots
  \label{eq:alphas}
\end{equation}
compared to the CODATA~2022 value $\alpha_s(M_Z) = 0.1179(9)$
($u_r = 7.6\times 10^{-3}$). The deviation is
$\delta = (0.117894-0.1179)/0.1179 = -48\ppm$. Since
$|\delta|=48\ppm \ll u_r = 7600\ppm$, the formula is well within current
experimental bounds.

\paragraph{Proton-to-electron mass ratio.}
\begin{equation}
  \frac{m_p}{m_e} = 6\pi^{5} = 6\times 305.020\ldots = 1836.118\ldots
  \label{eq:mpme}
\end{equation}
The CODATA~2022 value is $\mu = 1836.152\,673\,426(32)$
($u_r = 1.7\times 10^{-11}$). The deviation is
$\delta = -19\ppm$. Since $u_r = 17\ppt$, the formula deviates by
$|\delta|/u_r \approx 10^{6}$ times the experimental uncertainty. This formula
is therefore \textbf{already falsified as exact} by current measurements, but
it remains a striking approximation accurate to~5~significant figures.

\subsection{Error distribution analysis}\label{sec:errors}

The 16~approximate formulas (excluding the two exact results and the Pellis
formula) exhibit a characteristic error distribution:

\begin{itemize}
\item \textbf{4 formulas} with $|\delta|<50\ppm$: $m_p/m_e$, $\alpha_s$,
  $\sin^{2}\!\theta_{23}$, $\delta_{\mathrm{CP}}$.
\item \textbf{6 formulas} with $50\leq|\delta|<100\ppm$: $M_Z$,
  $\sin^{2}\!\theta_{12}$, $\sin^{2}\!\theta_{13}$, $T_{\mathrm{CMB}}$,
  $\sin^{2}\!\theta_W$.
\item \textbf{4 formulas} with $100\leq|\delta|<600\ppm$: $M_W$, $M_H$,
  $V_{us}$.
\item \textbf{1 formula} with $|\delta|\approx 900\ppm$: $G$.
\end{itemize}

The median deviation is approximately $85\ppm$. Notably, the electroweak
boson masses ($M_Z$, $M_W$) and the Higgs mass ($M_H$) all appear in the
100--200\ppm range, suggesting that the monomial framework captures the
leading-order structure of electroweak symmetry breaking but misses
higher-order corrections.

\subsection{Precision versus breadth}\label{sec:precision}

The two frameworks occupy complementary positions in the precision--breadth
trade-off:

\begin{itemize}
\item \textbf{Pellis (polynomial):} Achieves $\delta = -0.09\ppb$ for a
single constant ($\alpha^{-1}$) using three terms. The polynomial structure
allows each successive term to ``correct'' the previous approximation,
analogous to a perturbative series.

\item \textbf{Trinity (monomial):} Achieves $|\delta| < 100\ppm$ for
$\sim 10$~constants simultaneously using a single-term template. The monomial
constraint sacrifices per-formula precision for a unified algebraic structure.
\end{itemize}

This trade-off is not inherently an advantage of either approach; it reflects
different design philosophies. Polynomial frameworks are naturally suited to
high-precision approximation of individual constants, while monomial
frameworks are suited to identifying common mathematical structure across
many constants.

% ============================================================================
\section{Discussion}\label{sec:discussion}

\subsection{Structural complementarity: polynomial interference vs. monomial scaling}

The deepest structural observation of this comparison is the complementarity
between the Pellis and Trinity approaches. In the Pellis framework,
$\alpha^{-1}$ emerges as a sum of three terms that partially cancel
(interference): $137.508 - 0.472 + 0.000291 = 137.036\ldots$. In the Trinity
framework, the strong coupling emerges as a pure product (scaling):
$4\times 2.618 / (9\times 9.870) = 0.1179\ldots$.

This mirrors a well-known pattern in physics: Fourier series (polynomial,
interference) vs. power laws (monomial, scaling) often describe the same
phenomenon at different levels of approximation. It is natural to ask whether
the Trinity monomials are ``coarse-grained'' limits of underlying polynomial
structures---whether each monomial formula is the leading term of a more
detailed Pellis-type expansion.

\subsection{Renormalization hypothesis}

A speculative but testable hypothesis is that the monomial formulas correspond
to renormalization-group fixed points, while the polynomial corrections capture
the running between scales. In the renormalization group, coupling constants at
a fixed point often take simple, scale-invariant forms (monomials in fundamental
parameters), while their values at finite energy involve logarithmic corrections
(which could manifest as additive polynomial terms in $\vp^{-n}$).

This hypothesis would predict that: (a) constants known at high precision
($\alpha^{-1}$, $\mu$) should admit polynomial improvements over their monomial
approximations, and (b) constants known only at low precision ($G$,
$\alpha_s$) should currently be indistinguishable from their monomial
forms. The existing data are consistent with this prediction but do not
constitute evidence for it.

\subsection{Can $\pi$ and $e$ emerge from $\vp$-recursive processes?}

A deeper question is whether the mathematical constants $\pi$ and $e$
appearing in the Sacred Formula are themselves derivable from
$\vp$-recursive processes. Several mathematical identities connect these
constants:
\begin{align}
  \pi &= 4\sum_{k=0}^{\infty}\frac{(-1)^{k}}{2k+1}, &
  e &= \sum_{k=0}^{\infty}\frac{1}{k!},
\end{align}
and both arise naturally in the theory of continued fractions, where $\vp$
plays a distinguished role as the ``most irrational'' number (slowest
convergence of continued fraction approximants). If a single
$\vp$-recursive dynamical system could generate $\pi$, $e$, and $\gamma$
as natural outputs, the Sacred Formula would reduce to a structure with a
single free parameter~($\vp$).

\subsection{The $G$ problem}

The gravitational constant~$G$ presents a unique challenge. Its CODATA~2022
relative uncertainty is $u_r = 2.2\times 10^{-5}$ (22{,}000\ppm)---three to
six orders of magnitude larger than other constants in our table. The
Trinity formula $G=\pi^{3}\gamma^{2}/\vp$ deviates by $\sim 900\ppm$, which
is well within experimental uncertainty. However, this means that
\emph{any} formula producing a value within 2\% of the measured~$G$ would
be equally ``consistent'' with experiment. The $G$ formula cannot be
meaningfully tested until gravitational measurements improve by at least two
orders of magnitude.

\subsection{What would constitute falsification}\label{sec:falsify}

An empirical formula $V_{\mathrm{formula}}$ is falsified when either:
\begin{enumerate}[label=(\alph*)]
\item A new CODATA adjustment shifts the central value such that the formula's
  deviation increases beyond a significance threshold, or
\item Improved measurements reduce the experimental uncertainty below
  $|\delta|$, proving that the formula's prediction is inconsistent with
  experiment.
\end{enumerate}

Several formulas in our table are already in tension with experiment by
criterion~(b):
\begin{itemize}
\item $m_p/m_e = 6\pi^{5}$ has $|\delta|=19\ppm$ but $u_r = 17\ppt$---the
  formula is excluded as exact at the $10^{6}\sigma$ level.
\item $M_Z = 7\pi^{4}\vp e^{3}/243$ has $|\delta|=61\ppm$ but
  $u_r = 23\ppm$---the formula deviates at the $\sim 2.7\sigma$ level.
\end{itemize}

Conversely, formulas such as $\alpha_s = 4\vp^{2}/(9\pi^{2})$
($|\delta|=48\ppm$, $u_r = 7600\ppm$) and $G=\pi^{3}\gamma^{2}/\vp$
($|\delta|\sim 900\ppm$, $u_r = 22{,}000\ppm$) remain well within
experimental bounds and cannot currently be distinguished from
coincidence.

The Pellis formula $\alpha^{-1} = 360\vp^{-2}-2\vp^{-3}+(3\vp)^{-5}$ is
the most stringently tested: its $|\delta|=0.089\ppb$ is within
$u_r = 0.16\ppb$, but future improvements in the electron anomalous magnetic
moment measurement (which determines~$\alpha$) may either confirm or exclude
it at the ppb level.

\subsection*{Epistemological note}

\begin{quote}
\textbf{Disclaimer:} The formulas presented in this paper are empirical
approximations. They reproduce CODATA~2022 values to varying degrees of
precision but do not constitute derivations from first principles. No
Lagrangian, renormalization-group equation, or symmetry argument from which
the Sacred Formula $V = n\cdot 3^{k}\pi^{m}\vp^{p}e^{q}\gamma^{r}$ follows
has been identified. Close numerical agreement does not imply causal or
structural necessity. The standard of proof in physics requires a predictive
framework that (a)~derives the formula from a deeper physical principle,
(b)~predicts new observables not used in fitting, and (c)~withstands improved
experimental precision. This paper makes no such claim and invites critical
evaluation.
\end{quote}

% ============================================================================
\section{Conclusion}\label{sec:conclusion}

We have systematically compared 18~empirical formulas from the Trinity monomial
framework and the Pellis polynomial framework against CODATA~2022 recommended
values for fundamental physical constants. The Trinity identity
$\vp^{2}+\vp^{-2}=3$ is algebraically exact. Sixteen additional formulas
achieve deviations below $1000\ppm$; ten fall below $100\ppm$. The Pellis
three-term formula for $\alpha^{-1}$ achieves $\delta = -0.09\ppb$, within the
current experimental uncertainty of $0.16\ppb$.

We emphasize that these results constitute \textbf{empirical approximations,
not physical derivations}: no Lagrangian or symmetry principle from which the
Sacred Formula follows has been identified. The structural complementarity
between polynomial and monomial representations---additive interference versus
multiplicative scaling---merits further mathematical investigation. Whether
these patterns reflect deep structure in the fundamental constants or are
elaborate numerical coincidences remains an open question that only improved
experimental precision and theoretical insight can resolve.

\medskip
\noindent\textbf{Data availability.} All Python verification scripts are
available at \url{https://github.com/trinity-project}.

\medskip
\noindent\textbf{Author contributions.}
Conceptualization, D.V.\ and S.P.; Methodology, D.V.\ and S.P.;
Software, D.V.; Formal Analysis, D.V.\ and S.P.;
Writing---original draft, D.V.; Writing---review and editing, S.P.;
Visualization, D.V.
All authors have read and agreed to the published version of the manuscript.

% ============================================================================
\appendix

\section{Full Formula Catalog}\label{app:catalog}

Table~\ref{tab:catalog} lists all 18~formulas analyzed in this work with
their symbolic forms, exponent decompositions, predicted values, CODATA~2022
values, and relative deviations.

\begin{table}[H]
\centering
\caption{Complete catalog of formulas analyzed in this work. Tier~0 formulas
  are algebraically exact. For the Pellis formula (P1), the polynomial
  structure does not conform to the monomial template.}
\label{tab:catalog}
\small
\begin{tabular}{@{}clp{3.5cm}lll@{}}
\toprule
\# & Tier & Formula & Predicted & CODATA~2022 & $\delta$ \\
\midrule
1  & 0 & $\Ngen=\vp^{2}+\vp^{-2}$ & 3 & 3 gen. & EXACT \\
2  & 0 & $\theta_{\mathrm{QCD}}=|\vp^{2}+\vp^{-2}-3|$ & 0 & $<10^{-10}$ & EXACT \\
3  & 1 & $\delta_{\mathrm{CP}}=8\pi^{3}/(9e^{2})$ & $3.72999$ rad & $3.73$ rad & $+1.6\ppm$ \\
4  & 1 & $m_p/m_e=6\pi^{5}$ & $1836.118$ & $1836.15267(3)$ & $-19\ppm$ \\
5  & 1 & $\sin^{2}\!\theta_{23}$ & $0.5460$ & $0.546$ & $+28\ppm$ \\
6  & 1 & $\alpha_s=4\vp^{2}/(9\pi^{2})$ & $0.11789$ & $0.1179(9)$ & $-48\ppm$ \\
7  & 1 & $M_Z=7\pi^{4}\vp e^{3}/243$ & $91.193$ GeV & $91.1876(21)$ & $+61\ppm$ \\
8  & 1 & $\sin^{2}\!\theta_{12}$ & $0.30702$ & $0.307$ & $+75\ppm$ \\
9  & 1 & $\sin^{2}\!\theta_{13}$ & $0.02200$ & $0.0220$ & $+76\ppm$ \\
10 & 1 & $\sin^{2}\!\theta_W=2\pi^{3}e/729$ & $0.23123$ & $0.23121(4)$ & $+92\ppm$ \\
11 & 1 & $T_{\mathrm{CMB}}=\frac{5\pi^{4}\vp^{5}}{729e}$ & $2.72575$ K & $2.72548(57)$ & $+98\ppm$ \\
12 & 2 & $M_W=162\vp^{3}/(\pi e)$ & $80.359$ GeV & $80.369(13)$ & $-127\ppm$ \\
13 & 2 & $M_H=135\vp^{4}/e^{2}$ & $125.23$ GeV & $125.20(11)$ & $+210\ppm$ \\
14 & 2 & $V_{us}=3\gamma/\pi$ & $0.22543$ & $0.22530$ & $+570\ppm$ \\
15 & 2 & $G=\pi^{3}\gamma^{2}/\vp$ & ${\sim}6.674$ & $6.67430(15)$ & ${\sim}900\ppm$ \\
\midrule
P1 & --- & $\alpha^{-1}=360\vp^{-2}-2\vp^{-3}+(3\vp)^{-5}$ & $137.035999165$ & $137.035999177(21)$ & $-0.09\ppb$ \\
\bottomrule
\end{tabular}
\end{table}

\section{Computational Details}\label{app:code}

The following Python script was used to verify all numerical values in this
paper. It uses the \texttt{decimal} module with 50-digit precision to avoid
floating-point rounding errors.

\begin{small}
\begin{verbatim}
from decimal import Decimal, getcontext
getcontext().prec = 50

sqrt5 = Decimal(5).sqrt()
phi = (1 + sqrt5) / 2
pi = Decimal('3.14159265358979323846264338327950288419716939937510')
e  = Decimal('2.71828182845904523536028747135266249775724709369995')
gamma = phi**(-3)

CODATA = {
    'alpha_inv':    Decimal('137.035999177'),
    'mu':           Decimal('1836.152673426'),
    'alpha_s':      Decimal('0.1179'),
    'sin2_thetaW':  Decimal('0.23121'),
    'T_CMB':        Decimal('2.72548'),
    'M_H':          Decimal('125.20'),
    'M_Z':          Decimal('91.1876'),
    'M_W':          Decimal('80.369'),
}

def delta_ppm(formula_val, codata_val):
    return float((formula_val - codata_val) / codata_val) * 1e6

# Trinity identity
print(f"phi^2 + phi^-2 = {phi**2 + phi**(-2)}")  # 3.000...

# Pellis alpha^-1
pellis = 360*phi**(-2) - 2*phi**(-3) + (3*phi)**(-5)
print(f"Pellis alpha^-1 = {pellis}")
print(f"delta = {delta_ppm(pellis, CODATA['alpha_inv'])*1e3:.3f} ppb")

# Trinity formulas
formulas = {
    'mp/me = 6*pi^5':          (6*pi**5,           'mu'),
    'alpha_s = 4*phi^2/(9*pi^2)': (4*phi**2/(9*pi**2), 'alpha_s'),
    'sin2_thetaW':             (2*pi**3*e/729,      'sin2_thetaW'),
    'T_CMB':                   (5*pi**4*phi**5/(729*e), 'T_CMB'),
    'M_Z':                     (7*pi**4*phi*e**3/243, 'M_Z'),
    'M_H':                     (135*phi**4/e**2,     'M_H'),
    'M_W':                     (162*phi**3/(pi*e),   'M_W'),
}
for name, (val, key) in formulas.items():
    d = delta_ppm(val, CODATA[key])
    print(f"{name}: {float(val):.10f}, delta = {d:.1f} ppm")
\end{verbatim}
\end{small}

\section{CODATA 2022 Source Data}\label{app:codata}

Table~\ref{tab:codata} lists all CODATA~2022 values used in this work. All
values are from the NIST CODATA~2022 database~\cite{codata2022}.

\begin{table}[H]
\centering
\caption{CODATA~2022 recommended values used in this work.
  Source: NIST~\cite{codata2022},
  \url{https://physics.nist.gov/cuu/Constants/}.}
\label{tab:codata}
\begin{tabular}{llll}
\toprule
Constant & Symbol & CODATA~2022 Value & Rel.\ Uncertainty \\
\midrule
Fine-structure constant (inv.) & $\alpha^{-1}$ & $137.035\,999\,177(21)$ & $1.6\times 10^{-10}$ \\
Proton/electron mass ratio & $\mu=m_p/m_e$ & $1836.152\,673\,426(32)$ & $1.7\times 10^{-11}$ \\
Gravitational constant & $G$ & $6.67430(15)\times 10^{-11}$ m$^{3}$\,kg$^{-1}$\,s$^{-2}$ & $2.2\times 10^{-5}$ \\
Strong coupling & $\alpha_s(M_Z)$ & $0.1179(9)$ & $7.6\times 10^{-3}$ \\
Weak mixing angle & $\sin^{2}\!\theta_W$ & $0.23121(4)$ & $1.7\times 10^{-4}$ \\
CMB temperature & $T_{\mathrm{CMB}}$ & $2.72548(57)$ K & $2.1\times 10^{-4}$ \\
Higgs boson mass & $M_H$ & $125.20(11)$ GeV & $8.8\times 10^{-4}$ \\
Z boson mass & $M_Z$ & $91.1876(21)$ GeV & $2.3\times 10^{-5}$ \\
W boson mass & $M_W$ & $80.369(13)$ GeV & $1.6\times 10^{-4}$ \\
\bottomrule
\end{tabular}
\end{table}

% ============================================================================
\begin{thebibliography}{99}

\bibitem{feynman1985}
R.\,P.~Feynman,
\textit{QED: The Strange Theory of Light and Matter},
Princeton University Press, 1985.

\bibitem{codata2022}
{NIST/CODATA},
``{CODATA Internationally Recommended 2022 Values of the Fundamental Physical Constants},''
\url{https://physics.nist.gov/cuu/Constants/}, 2022.

\bibitem{pellis2021alpha}
S.~Pellis,
``Exact Formula for the Fine-Structure Constant $\alpha$ in Terms of the Golden Ratio $\varphi$,''
viXra:2110.0053, 2021.
\url{https://vixra.org/abs/2110.0053}

\bibitem{pellis2023toe}
S.~Pellis,
``Dimensionless Theory of Everything,''
SSRN 4469668, 2023.
\url{https://papers.ssrn.com/sol3/papers.cfm?abstract_id=4469668}

\bibitem{heyrovska2005}
R.~Heyrovska,
``Fine-structure Constant, Anomalous Magnetic Moment, Relativity Factor and the Golden Ratio that Divides the Bohr Radius,''
arXiv:physics/0509207, 2005.

\bibitem{coldea2010}
R.~Coldea \textit{et al.},
``Quantum Criticality in an Ising Chain: Experimental Evidence for Emergent E$_{8}$ Symmetry,''
\textit{Science} \textbf{327}, 177--180 (2010).
DOI: 10.1126/science.1180085.

\bibitem{sherbon2018}
M.\,A.~Sherbon,
``Golden Ratio Geometry and the Fine-Structure Constant,''
SSRN 3464431, 2018.
\url{https://papers.ssrn.com/sol3/papers.cfm?abstract_id=3464431}

\bibitem{sommerfeld1916}
A.~Sommerfeld,
``Zur Quantentheorie der Spektrallinien,''
\textit{Ann.\ Phys.} \textbf{51}, 1--94 (1916).

\end{thebibliography}

\end{document}
